% Body of supplement-acm.tex (electronic supplement), generated by
% pddl/sync_from_paper.py. Do not edit; edit the paper and re-sync.
\section{Algorithms of the \plansys extensions}\label{planta:S-sec:plansys_algorithms}

This section details the algorithms of the \plansys extensions described in \Cref{planta:sec:impl_plansys}.

To support the evaluation of existential preconditions and derived predicates in \plansys, this work proposes the evaluation algorithm described in \Cref{planta:S-alg:eval}. Lines~2-3 evaluate predicates by unifying their free parameters with all objects in $\pddlState$ that satisfy the predicate,  returning a set of dictionaries of the form $\{\{?x_1:o_1, \ldots{},?x_n:o_n\} \mid (p~o_1,\ldots{},o_n) \in S\}$. Lines 4-6 solve negated conditions by evaluating the condition and returning the complement of the evaluated set $E$ with respect to the current state $\pddlState$, with the form $\{\{?x_1:o_1,\ldots{},?x_n:o_n\} \mid (o_1, \ldots{}, o_n)  \in S \wedge  \{?x_1:o_1, \ldots{},?x_n:o_n\} \notin E\}$. Lines 7-13 evaluate conjunctions of conditions by individually evaluating each condition and joining the solution sets. Lines 14-17 evaluate disjunctions by evaluating each condition and merging the solution sets. Lines 18-20 evaluate existential preconditions by evaluating their body and returning a dictionary without free variables in the form $\{\{?x_1:o_1,\ldots{},?x_n:o_n\} \mid \{?x_1:o_1,\ldots{},?x_n:o_n,?x_{e_{1}}:o_{e_{1}},\ldots{}, ?x_{e_{k}}:o_{e_{k}}\} \in E\}$.

\begin{algorithm}[tb]
    \caption{Evaluation algorithm}%
    \label{planta:S-alg:eval}
    \begin{algorithmic}[1] % The number tells where the line numbering should start
        \Procedure{Evaluate}{$\algC, \algS$}
            \If{$ \algC = (p~?x_1 \ldots{} ?x_n)$}
                \State \textbf{return} $\Call{UnifyPredicate}{\algC, \algS}$
            \ElsIf{$\algC = (not~(\algC'))$}
                \State $\algE \gets \Call{Evaluate}{\algC', \algS}$
                \State \textbf{return} $\Call{ComplementSet}{\algE, \algS}$
            \ElsIf{$\algC = (and~(\algC'~\ldots{}~\algC''))$}
                \State $\algE \gets \{\emptyset\}$
                \For{$c \in \{\algC', \ldots{}, \algC''\}$}
                    \State $\algE' \gets \Call{Evaluate}{c,~\algS}$
                    \State $\algE \gets \Call{Join}{\algE,~\algE'}$
                \EndFor
                \State \textbf{return} $\algE$
            \ElsIf{$\algC = (or~(\algC'~\algC''))$}
                \State $\algE' \gets \Call{Evaluate}{\algC',~\algS}$
                \State $\algE'' \gets \Call{Evaluate}{\algC'',~\algS}$
                \State \textbf{return} $\Call{Union}{\algE',~\algE''}$
             \ElsIf{$\algC = (exists~(?x_{e_{1}} \ldots{} ?x_{e_{k}})~\algC')$}
                \State $\algE \gets \Call{Evaluate}{\algC',~\algS}$
                \State \textbf{return} $\algE \setminus \{?x_{e_{1}},\ldots{},?x_{e_{k}}\}$
            \EndIf
        \EndProcedure
    \end{algorithmic}
\end{algorithm}

To solve derived predicates, this work introduces Algorithm~\ref{planta:S-alg:solve_derived}. It partitions the resolution graph into acyclic and cyclic components. The acyclic components are resolved in topological order, while the cyclic components are resolved using a fixed-point iteration. Root nodes are considered solved by definition, as they correspond to non-derived predicates. For instance, in the graph shown in \Cref{planta:fig:derived_graph}, the derived predicates $\texttt{inferredB}$ and $\texttt{inferredAB}$ form a cycle, whereas the other derived predicates are acyclic.

Line~2 calls the $\texttt{GetDerivedSCCs}$ function, which computes the strongly connected components~(SCCs) of the subgraph formed by traversing all descendant nodes of the $\texttt{Roots}$ nodes in the resolution graph, using Tarjan's algorithm~\cite{doi:10.1137/0201010}. The function returns a vector of SCCs in topological order. Each SCC contains only derived predicates, explicitly excluding predicates, functions, and instances. For example, if $\texttt{Roots} = \{\texttt{predicateA}, \texttt{predicateB}\}$,
\begin{equation*}
\texttt{GetDerivedSCCs}(\texttt{Graph}, \texttt{Roots}) = \big\{ \{\texttt{inferredA}\},\, \{\texttt{inferredAA}\},\, \{\texttt{inferredB},\, \texttt{inferredAB}\},\, \{\texttt{inferredBB}\} \big\}.
\end{equation*}

Lines~4--15 solve the cyclic SCCs using a fixed-point iteration, while lines~17--22 handle each acyclic SCC by solving its single derived predicate. The $\texttt{Unground}$ function removes all existing assertions of $d$ and its descendants from the state $\pddlState$. The $\texttt{Evaluate}$ function evaluates the preconditions of the derived predicate $d$, storing all valid parameter bindings in the set $E$ (e.g., $\{?x = \texttt{instance1}\}$). The $\texttt{GroundPredicates}$ function grounds the derived predicate for each substitution in $E$ (e.g., $(\texttt{inferredA}~\texttt{instance1})$) and adds the resulting assertions to the state.

\begin{algorithm}[tb]
\caption{Derived Predicates Evaluation}%
\label{planta:S-alg:solve_derived}
    \begin{algorithmic}[1]
    \Procedure{SolveDerived}{$\algGraph, \algRoots, \algS$}
        \State $\algSCCs \gets \Call{GetDerivedSCCs}{\algGraph, \algRoots}$
        \For{$\algSCC \in \algSCCs$}
            \If{$\Call{IsCyclic}{\algSCC}$}
                \For{$d \in \algSCC$}
                    \State $\algS \gets \Call{Unground}{\algS, d}$
                \EndFor
                \Repeat
                    \State $\algS' \gets \algS$
                    \For{$d \in \algSCC$}
                        \State $\algE \gets \Call{Evaluate}{\mathrm{pre}(d), \algS}$
                        \State $\algP' \gets \Call{GroundPredicates}{d, \algE}$
                        \State $\algS \gets \algS \cup \algP'$
                    \EndFor
                \Until{$\algS = \algS'$}
            \Else
                \For{$d \in \algSCC$}
                    \State $\algS \gets \Call{Unground}{\algS, d}$
                    \State $\algE \gets \Call{Evaluate}{\mathrm{pre}(d), \algS}$
                    \State $\algP' \gets \Call{GroundPredicates}{d, \algE}$
                    \State $\algS \gets \algS \cup \algP'$
                \EndFor
            \EndIf
        \EndFor
    \State \Return $\algS$
    \EndProcedure
    \end{algorithmic}
\end{algorithm}

\Cref{planta:S-alg:apply_new} defines the $\texttt{ApplyEffects}$ function. Its inputs are the current state $\algS$, the effects $\algEff$ of an action, the resolution graph $\algGraph$, the set $\algNodes$ of modified nodes, and a boolean $\algDerive$ that indicates whether derived predicates should be recomputed. Lines~2--4 add a predicate to the state, and lines~5--7 remove a negated predicate from it. In both cases the predicate is added to $\algNodes$. Lines~8--11 apply conjunctive effects recursively, with $\algDerive$ set to false so that derived predicates are not recomputed after each individual effect. If $\algDerive$ is true (line~13), the function calls $\texttt{SolveDerived}$ (\Cref{planta:S-alg:solve_derived}) with $\algNodes$ as root nodes, so only the derived predicates that depend on the modified predicates are re-evaluated. The function returns the updated state and the set of modified nodes.

\begin{algorithm}[tb]
    \caption{Apply Action Effects and Selectively Re-evaluate Derived Predicates}%
    \label{planta:S-alg:apply_new}
    \begin{algorithmic}[1]
    \Procedure{ApplyEffects}{$\algS, \algEff, \algGraph, \algNodes, \algDerive$}
        \If{$\algEff = (p~\underline{o}_1~\dots~\underline{o}_n)$}
            \State $\algS \gets \algS \cup \{\algEff\}$
            \State $\algNodes \gets \algNodes \cup \{\algEff\}$
        \ElsIf{$\algEff = (not~(\algEff'))$}
            \State $\algS \gets \algS \setminus \{\algEff'\}$
            \State $\algNodes \gets \algNodes \cup \{\algEff'\}$
        \ElsIf{$\algEff = (and~(\algEff'_1~\dots~\algEff'_k))$}
            \For{$e \in \{\algEff'_1, \dots, \algEff'_k\}$}
                \State $\algS, \algNodes \gets \Call{ApplyEffects}{\algS, e, \algGraph, \algNodes, \text{false}}$
            \EndFor
        \EndIf
        \If{$\algDerive$}
            \State $\algS \gets \Call{SolveDerived}{\algGraph, \algNodes, \algS}$
        \EndIf
        \State \Return $\algS, \algNodes$
    \EndProcedure
    \end{algorithmic}
\end{algorithm}

\section{Wilcoxon signed-rank comparisons for the SUAVE experiments}\label{planta:S-sec:wilcoxon}

This section reports the paired Wilcoxon signed-rank comparisons for experiments 1 and 2 and extended experiments 1 and 2 of the SUAVE exemplar. The test setup and hypotheses are described in \Cref{planta:sec:suave_exp}, and the comparisons for experiment 3 and extended experiment 3 are reported in \Cref{planta:tab:wilcoxon}, whose caption also explains the layout and color legend of the tables below.

\begin{table}[tb]
\centering
\caption{Experiment 1: paired Wilcoxon signed-rank comparisons.}%
\label{planta:S-tab:wilcoxon-exp1}
\inputtablebody{pddl/tables/wilcoxon_analysis/exp1_wilcoxon.tex}
\end{table}

\begin{table}[tb]
\centering
\caption{Experiment 2: paired Wilcoxon signed-rank comparisons.}%
\label{planta:S-tab:wilcoxon-exp2}
\inputtablebody{pddl/tables/wilcoxon_analysis/exp2_wilcoxon.tex}
\end{table}

\begin{table}[tb]
\centering
\caption{Extended experiment 1: paired Wilcoxon signed-rank comparisons.}%
\label{planta:S-tab:wilcoxon-extended-exp1}
\inputtablebody{pddl/tables/wilcoxon_analysis/extended_exp1_wilcoxon.tex}
\end{table}

\begin{table}[tb]
\centering
\caption{Extended experiment 2: paired Wilcoxon signed-rank comparisons.}%
\label{planta:S-tab:wilcoxon-extended-exp2}
\inputtablebody{pddl/tables/wilcoxon_analysis/extended_exp2_wilcoxon.tex}
\end{table}

\section{Per-experiment results of the scalability experiments}\label{planta:S-sec:scalability_plots}

This section reports planning time and plan length for each scalability experiment of the navigation use case. The experimental design and setup are described in \Cref{planta:sec:scale_exp}, and the results are summarized in \Cref{planta:tab:scalability_summary} and \Cref{planta:fig:scalability_time}. In \Cref{planta:S-fig:s1_results,planta:S-fig:s4_results}, the no-adaptation curve reports A* search.

\begin{figure}[tb]
    \centering
    \begin{subfigure}[t]{0.44\linewidth}
        \centering
        % Generated by navigation_planta; do not edit by hand.
\begin{tikzpicture}[baseline=(paperAxis.north)]
\begin{axis}[name=paperAxis,
width=\textwidth,
height=0.75\textwidth,
xlabel={Number of Nodes},
ylabel={Planning Time [seconds]},
xmin=0,
xmax=900,
ymin=0,
ymax=1,
xtick={0,200,400,600,800,900},
grid=major,
legend pos=north west,
legend style={font=\footnotesize,nodes={scale=0.8,transform shape}},
legend image post style={scale=0.5}
]
\addplot[name path=upper0,draw=none,forget plot] table[x=nodes,y=mean_time_plus_std,col sep=comma] {pddl/navigation_experiment/data/s1_a_star_experiment_processed.csv};
\addplot[name path=lower0,draw=none,forget plot] table[x=nodes,y=mean_time_minus_std,col sep=comma] {pddl/navigation_experiment/data/s1_a_star_experiment_processed.csv};
\addplot[fill=\greenColor!20,fill opacity=0.3,draw=none,forget plot] fill between[of=upper0 and lower0];
\addplot[thick,color=\greenColor] table[x=nodes,y=mean_time,col sep=comma] {pddl/navigation_experiment/data/s1_a_star_experiment_processed.csv};
\addlegendentry{A* mean}
\addplot[name path=upper1,draw=none,forget plot] table[x=nodes,y=mean_time_plus_std,col sep=comma] {pddl/navigation_experiment/data/s1_eager_experiment_processed.csv};
\addplot[name path=lower1,draw=none,forget plot] table[x=nodes,y=mean_time_minus_std,col sep=comma] {pddl/navigation_experiment/data/s1_eager_experiment_processed.csv};
\addplot[fill=\blueColor!20,fill opacity=0.3,draw=none,forget plot] fill between[of=upper1 and lower1];
\addplot[thick,color=\blueColor] table[x=nodes,y=mean_time,col sep=comma] {pddl/navigation_experiment/data/s1_eager_experiment_processed.csv};
\addlegendentry{EG mean}
\addplot[name path=upper2,draw=none,forget plot] table[x=nodes,y=mean_time_plus_std,col sep=comma] {pddl/navigation_experiment/data/s1_a_star_no_sas_experiment_processed.csv};
\addplot[name path=lower2,draw=none,forget plot] table[x=nodes,y=mean_time_minus_std,col sep=comma] {pddl/navigation_experiment/data/s1_a_star_no_sas_experiment_processed.csv};
\addplot[fill=\orangeColor!20,fill opacity=0.3,draw=none,forget plot] fill between[of=upper2 and lower2];
\addplot[thick,color=\orangeColor] table[x=nodes,y=mean_time,col sep=comma] {pddl/navigation_experiment/data/s1_a_star_no_sas_experiment_processed.csv};
\addlegendentry{No adap. mean}
\end{axis}
\end{tikzpicture}

        \caption{Planning time as map size increases.}%
        \label{planta:S-plot:s1_planning_time}
    \end{subfigure}
    \begin{subfigure}[t]{0.44\linewidth}
        \centering
        % Generated by navigation_planta; do not edit by hand.
\begin{tikzpicture}[baseline=(paperAxis.north)]
\begin{axis}[name=paperAxis,
width=\textwidth,
height=0.75\textwidth,
xlabel={Number of Nodes},
ylabel={Plan Length [no. of actions]},
xmin=0,
xmax=900,
ymin=0,
ymax=70,
xtick={0,200,400,600,800,900},
grid=major,
legend pos=north west,
legend style={font=\footnotesize,nodes={scale=0.8,transform shape}},
legend image post style={scale=0.5}
]
\addplot[name path=upper0,draw=none,forget plot] table[x=nodes,y=mean_action_count_plus_std,col sep=comma] {pddl/navigation_experiment/data/s1_a_star_experiment_processed.csv};
\addplot[name path=lower0,draw=none,forget plot] table[x=nodes,y=mean_action_count_minus_std,col sep=comma] {pddl/navigation_experiment/data/s1_a_star_experiment_processed.csv};
\addplot[fill=\greenColor!20,fill opacity=0.3,draw=none,forget plot] fill between[of=upper0 and lower0];
\addplot[thick,color=\greenColor] table[x=nodes,y=mean_action_count,col sep=comma] {pddl/navigation_experiment/data/s1_a_star_experiment_processed.csv};
\addlegendentry{A* mean}
\addplot[name path=upper1,draw=none,forget plot] table[x=nodes,y=mean_action_count_plus_std,col sep=comma] {pddl/navigation_experiment/data/s1_eager_experiment_processed.csv};
\addplot[name path=lower1,draw=none,forget plot] table[x=nodes,y=mean_action_count_minus_std,col sep=comma] {pddl/navigation_experiment/data/s1_eager_experiment_processed.csv};
\addplot[fill=\blueColor!20,fill opacity=0.3,draw=none,forget plot] fill between[of=upper1 and lower1];
\addplot[thick,color=\blueColor] table[x=nodes,y=mean_action_count,col sep=comma] {pddl/navigation_experiment/data/s1_eager_experiment_processed.csv};
\addlegendentry{EG mean}
\end{axis}
\end{tikzpicture}

        \caption{Plan length as map size increases. Both search strategies overlap.}%
        \label{planta:S-plot:s1_action_number}
    \end{subfigure}
    \caption{\Sone results: Planning time and plan length as map size increases. The shaded area shows the standard deviation. EG stands for the eager greedy search.}%
    \label{planta:S-fig:s1_results}
\Description{Two line-plot panels show mean planning time in seconds and mean plan length in action
counts against map size, from 9 to 900 retained nodes. A* and eager greedy search have
similar planning times, reaching 0.92 and 0.95 seconds, respectively, at 900 nodes, and overlapping
plan-length curves, reaching 47.20 actions. The left panel also shows the faster
no-adaptation baseline. Shaded bands indicate standard deviations.}
\end{figure}

\begin{figure}[tb]
    \centering
    \begin{subfigure}[t]{0.44\linewidth}
        \centering
        % Generated by navigation_planta; do not edit by hand.
\begin{tikzpicture}[baseline=(paperAxis.north)]
\begin{axis}[name=paperAxis,
width=\textwidth,
height=0.75\textwidth,
xlabel={Number of arch. configurations},
ylabel={Planning Time [seconds]},
xmin=2,
xmax=15,
ymin=0,
ymax=0.32,
xtick={2,4,6,8,10,12,14},
grid=major,
legend pos=north west,
legend style={font=\footnotesize,nodes={scale=0.8,transform shape}},
legend image post style={scale=0.5}
]
\addplot[name path=upper0,draw=none,forget plot] table[x=n_fd,y=mean_time_plus_std,col sep=comma] {pddl/navigation_experiment/data/s2_astar_experiment_processed.csv};
\addplot[name path=lower0,draw=none,forget plot] table[x=n_fd,y=mean_time_minus_std,col sep=comma] {pddl/navigation_experiment/data/s2_astar_experiment_processed.csv};
\addplot[fill=\greenColor!20,fill opacity=0.3,draw=none,forget plot] fill between[of=upper0 and lower0];
\addplot[thick,color=\greenColor] table[x=n_fd,y=mean_time,col sep=comma] {pddl/navigation_experiment/data/s2_astar_experiment_processed.csv};
\addlegendentry{A* mean}
\addplot[name path=upper1,draw=none,forget plot] table[x=n_fd,y=mean_time_plus_std,col sep=comma] {pddl/navigation_experiment/data/s2_eager_experiment_processed.csv};
\addplot[name path=lower1,draw=none,forget plot] table[x=n_fd,y=mean_time_minus_std,col sep=comma] {pddl/navigation_experiment/data/s2_eager_experiment_processed.csv};
\addplot[fill=\blueColor!20,fill opacity=0.3,draw=none,forget plot] fill between[of=upper1 and lower1];
\addplot[thick,color=\blueColor] table[x=n_fd,y=mean_time,col sep=comma] {pddl/navigation_experiment/data/s2_eager_experiment_processed.csv};
\addlegendentry{EG mean}
\end{axis}
\end{tikzpicture}

        \caption{Planning time as architectural configurations increase.}%
        \label{planta:S-plot:s2_planning_time}
    \end{subfigure}
    \begin{subfigure}[t]{0.44\linewidth}
        \centering
        % Generated by navigation_planta; do not edit by hand.
\begin{tikzpicture}[baseline=(paperAxis.north)]
\begin{axis}[name=paperAxis,
width=\textwidth,
height=0.75\textwidth,
xlabel={Number of arch. configurations},
ylabel={Plan Length [no. of actions]},
xmin=2,
xmax=15,
ymin=0,
ymax=32,
xtick={2,4,6,8,10,12,14},
grid=major,
legend pos=north west,
legend style={font=\footnotesize,nodes={scale=0.8,transform shape}},
legend image post style={scale=0.5}
]
\addplot[name path=upper0,draw=none,forget plot] table[x=n_fd,y=mean_action_count_plus_std,col sep=comma] {pddl/navigation_experiment/data/s2_astar_experiment_processed.csv};
\addplot[name path=lower0,draw=none,forget plot] table[x=n_fd,y=mean_action_count_minus_std,col sep=comma] {pddl/navigation_experiment/data/s2_astar_experiment_processed.csv};
\addplot[fill=\greenColor!20,fill opacity=0.3,draw=none,forget plot] fill between[of=upper0 and lower0];
\addplot[thick,color=\greenColor] table[x=n_fd,y=mean_action_count,col sep=comma] {pddl/navigation_experiment/data/s2_astar_experiment_processed.csv};
\addlegendentry{A* mean}
\addplot[name path=upper1,draw=none,forget plot] table[x=n_fd,y=mean_action_count_plus_std,col sep=comma] {pddl/navigation_experiment/data/s2_eager_experiment_processed.csv};
\addplot[name path=lower1,draw=none,forget plot] table[x=n_fd,y=mean_action_count_minus_std,col sep=comma] {pddl/navigation_experiment/data/s2_eager_experiment_processed.csv};
\addplot[fill=\blueColor!20,fill opacity=0.3,draw=none,forget plot] fill between[of=upper1 and lower1];
\addplot[thick,color=\blueColor] table[x=n_fd,y=mean_action_count,col sep=comma] {pddl/navigation_experiment/data/s2_eager_experiment_processed.csv};
\addlegendentry{EG mean}
\end{axis}
\end{tikzpicture}

        \caption{Plan length as architectural configurations increase. Both search strategies overlap.}%
        \label{planta:S-plot:s2_action_number}
    \end{subfigure}
    \caption{\Stwo results: Planning time and plan length as architectural configurations increase. The shaded area shows the standard deviation. EG stands for the eager greedy search.}%
    \label{planta:S-fig:s2_results}
\Description{Two line-plot panels show mean planning time in seconds and mean plan length in action
counts as the number of architectural configurations increases from 2 to 15. Planning
time changes modestly, reaching 0.30 seconds for both A* and eager
greedy search. The plan-length curves overlap, reaching 20.70 actions at 15
configurations. Shaded bands indicate standard deviations.}
\end{figure}

\begin{figure}[tb]
    \centering
    \begin{subfigure}[t]{0.44\linewidth}
        \centering
        % Generated by navigation_planta; do not edit by hand.
\begin{tikzpicture}[baseline=(paperAxis.north)]
\begin{axis}[name=paperAxis,
width=\textwidth,
height=0.75\textwidth,
xlabel={Number of corridor types},
ylabel={Planning Time [seconds]},
xmin=2,
xmax=4,
ymin=0,
ymax=0.32,
xtick={2,3,4},
grid=major,
legend pos=north west,
legend style={font=\footnotesize,nodes={scale=0.8,transform shape}},
legend image post style={scale=0.5}
]
\addplot[name path=upper0,draw=none,forget plot] table[x=k,y=mean_time_plus_std,col sep=comma] {pddl/navigation_experiment/data/s3_astar_experiment_processed.csv};
\addplot[name path=lower0,draw=none,forget plot] table[x=k,y=mean_time_minus_std,col sep=comma] {pddl/navigation_experiment/data/s3_astar_experiment_processed.csv};
\addplot[fill=\greenColor!20,fill opacity=0.3,draw=none,forget plot] fill between[of=upper0 and lower0];
\addplot[thick,color=\greenColor] table[x=k,y=mean_time,col sep=comma] {pddl/navigation_experiment/data/s3_astar_experiment_processed.csv};
\addlegendentry{A* mean}
\addplot[name path=upper1,draw=none,forget plot] table[x=k,y=mean_time_plus_std,col sep=comma] {pddl/navigation_experiment/data/s3_eager_experiment_processed.csv};
\addplot[name path=lower1,draw=none,forget plot] table[x=k,y=mean_time_minus_std,col sep=comma] {pddl/navigation_experiment/data/s3_eager_experiment_processed.csv};
\addplot[fill=\blueColor!20,fill opacity=0.3,draw=none,forget plot] fill between[of=upper1 and lower1];
\addplot[thick,color=\blueColor] table[x=k,y=mean_time,col sep=comma] {pddl/navigation_experiment/data/s3_eager_experiment_processed.csv};
\addlegendentry{EG mean}
\end{axis}
\end{tikzpicture}

        \caption{Planning time as corridor types increase.}%
        \label{planta:S-plot:s3_planning_time}
    \end{subfigure}
    \begin{subfigure}[t]{0.44\linewidth}
        \centering
        % Generated by navigation_planta; do not edit by hand.
\begin{tikzpicture}[baseline=(paperAxis.north)]
\begin{axis}[name=paperAxis,
width=\textwidth,
height=0.75\textwidth,
xlabel={Number of corridor types},
ylabel={Plan Length [no. of actions]},
xmin=2,
xmax=4,
ymin=0,
ymax=32,
xtick={2,3,4},
grid=major,
legend pos=north west,
legend style={font=\footnotesize,nodes={scale=0.8,transform shape}},
legend image post style={scale=0.5}
]
\addplot[name path=upper0,draw=none,forget plot] table[x=k,y=mean_action_count_plus_std,col sep=comma] {pddl/navigation_experiment/data/s3_astar_experiment_processed.csv};
\addplot[name path=lower0,draw=none,forget plot] table[x=k,y=mean_action_count_minus_std,col sep=comma] {pddl/navigation_experiment/data/s3_astar_experiment_processed.csv};
\addplot[fill=\greenColor!20,fill opacity=0.3,draw=none,forget plot] fill between[of=upper0 and lower0];
\addplot[thick,color=\greenColor] table[x=k,y=mean_action_count,col sep=comma] {pddl/navigation_experiment/data/s3_astar_experiment_processed.csv};
\addlegendentry{A* mean}
\addplot[name path=upper1,draw=none,forget plot] table[x=k,y=mean_action_count_plus_std,col sep=comma] {pddl/navigation_experiment/data/s3_eager_experiment_processed.csv};
\addplot[name path=lower1,draw=none,forget plot] table[x=k,y=mean_action_count_minus_std,col sep=comma] {pddl/navigation_experiment/data/s3_eager_experiment_processed.csv};
\addplot[fill=\blueColor!20,fill opacity=0.3,draw=none,forget plot] fill between[of=upper1 and lower1];
\addplot[thick,color=\blueColor] table[x=k,y=mean_action_count,col sep=comma] {pddl/navigation_experiment/data/s3_eager_experiment_processed.csv};
\addlegendentry{EG mean}
\end{axis}
\end{tikzpicture}

        \caption{Plan length as corridor types increase. Both search strategies overlap.}%
        \label{planta:S-plot:s3_action_number}
    \end{subfigure}
    \caption{\Sthree results: Planning time and plan length as corridor types increase. The shaded area shows the standard deviation. EG stands for the eager greedy search.}%
    \label{planta:S-fig:s3_results}
\Description{Two line-plot panels show mean planning time in seconds and mean plan length in action
counts as the number of corridor types increases from 2 to 4. A* and eager greedy search
have nearly identical planning times, increasing from about 0.15 to 0.20 seconds. Their
plan-length curves overlap, with means of 19.1 actions at two corridor types and 20.3 at
four. Shaded bands indicate standard deviations.}
\end{figure}

\begin{figure}[tb]
    \centering
    \begin{subfigure}[t]{0.44\linewidth}
        \centering
        % Generated by navigation_planta; do not edit by hand.
\begin{tikzpicture}[baseline=(paperAxis.north)]
\begin{axis}[name=paperAxis,
width=\textwidth,
height=0.75\textwidth,
xlabel={Number of mission actions},
ylabel={Planning Time [seconds]},
xmin=1,
xmax=5,
ymin=0,
ymax=18.3531757133,
xtick={1,2,3,4,5},
grid=major,
legend pos=north west,
legend style={font=\footnotesize,nodes={scale=0.8,transform shape}},
legend image post style={scale=0.5}
]
\addplot[name path=upper0,draw=none,forget plot] table[x=m,y=mean_time_plus_std,col sep=comma] {pddl/navigation_experiment/data/s4_astar_experiment_processed.csv};
\addplot[name path=lower0,draw=none,forget plot] table[x=m,y=mean_time_minus_std,col sep=comma] {pddl/navigation_experiment/data/s4_astar_experiment_processed.csv};
\addplot[fill=\greenColor!20,fill opacity=0.3,draw=none,forget plot] fill between[of=upper0 and lower0];
\addplot[thick,color=\greenColor] table[x=m,y=mean_time,col sep=comma] {pddl/navigation_experiment/data/s4_astar_experiment_processed.csv};
\addlegendentry{A* mean}
\addplot[name path=upper1,draw=none,forget plot] table[x=m,y=mean_time_plus_std,col sep=comma] {pddl/navigation_experiment/data/s4_eager_experiment_processed.csv};
\addplot[name path=lower1,draw=none,forget plot] table[x=m,y=mean_time_minus_std,col sep=comma] {pddl/navigation_experiment/data/s4_eager_experiment_processed.csv};
\addplot[fill=\blueColor!20,fill opacity=0.3,draw=none,forget plot] fill between[of=upper1 and lower1];
\addplot[thick,color=\blueColor] table[x=m,y=mean_time,col sep=comma] {pddl/navigation_experiment/data/s4_eager_experiment_processed.csv};
\addlegendentry{EG mean}
\addplot[name path=upper2,draw=none,forget plot] table[x=m,y=mean_time_plus_std,col sep=comma] {pddl/navigation_experiment/data/s4_astar_no_sas_experiment_processed.csv};
\addplot[name path=lower2,draw=none,forget plot] table[x=m,y=mean_time_minus_std,col sep=comma] {pddl/navigation_experiment/data/s4_astar_no_sas_experiment_processed.csv};
\addplot[fill=\orangeColor!20,fill opacity=0.3,draw=none,forget plot] fill between[of=upper2 and lower2];
\addplot[thick,color=\orangeColor] table[x=m,y=mean_time,col sep=comma] {pddl/navigation_experiment/data/s4_astar_no_sas_experiment_processed.csv};
\addlegendentry{No adap. mean}
\end{axis}
\end{tikzpicture}

        \caption{Planning time as mission-level actions increase.}%
        \label{planta:S-plot:s4_planning_time}
    \end{subfigure}
    \begin{subfigure}[t]{0.44\linewidth}
        \centering
        % Generated by navigation_planta; do not edit by hand.
\begin{tikzpicture}[baseline=(paperAxis.north)]
\begin{axis}[name=paperAxis,
width=\textwidth,
height=0.75\textwidth,
xlabel={Number of mission actions},
ylabel={Plan Length [no. of actions]},
xmin=1,
xmax=5,
ymin=0,
ymax=55,
xtick={1,2,3,4,5},
grid=major,
legend pos=north west,
legend style={font=\footnotesize,nodes={scale=0.8,transform shape}},
legend image post style={scale=0.5}
]
\addplot[name path=upper0,draw=none,forget plot] table[x=m,y=mean_action_count_plus_std,col sep=comma] {pddl/navigation_experiment/data/s4_astar_experiment_processed.csv};
\addplot[name path=lower0,draw=none,forget plot] table[x=m,y=mean_action_count_minus_std,col sep=comma] {pddl/navigation_experiment/data/s4_astar_experiment_processed.csv};
\addplot[fill=\greenColor!20,fill opacity=0.3,draw=none,forget plot] fill between[of=upper0 and lower0];
\addplot[thick,color=\greenColor] table[x=m,y=mean_action_count,col sep=comma] {pddl/navigation_experiment/data/s4_astar_experiment_processed.csv};
\addlegendentry{A* mean}
\addplot[name path=upper1,draw=none,forget plot] table[x=m,y=mean_action_count_plus_std,col sep=comma] {pddl/navigation_experiment/data/s4_eager_experiment_processed.csv};
\addplot[name path=lower1,draw=none,forget plot] table[x=m,y=mean_action_count_minus_std,col sep=comma] {pddl/navigation_experiment/data/s4_eager_experiment_processed.csv};
\addplot[fill=\blueColor!20,fill opacity=0.3,draw=none,forget plot] fill between[of=upper1 and lower1];
\addplot[thick,color=\blueColor] table[x=m,y=mean_action_count,col sep=comma] {pddl/navigation_experiment/data/s4_eager_experiment_processed.csv};
\addlegendentry{EG mean}
\end{axis}
\end{tikzpicture}

        \caption{Plan length as mission-level actions increase.}%
        \label{planta:S-plot:s4_action_number}
    \end{subfigure}
    \caption{\Sfour results: Planning time and plan length as mission-level actions increase. The shaded area shows the standard deviation. EG stands for the eager greedy search.}%
    \label{planta:S-fig:s4_results}
\Description{Two line-plot panels show mean planning time in seconds and mean plan length in action
counts as the number of mission-level action schemas increases from 1 to 5. A* planning
time rises sharply to 15.45 seconds, while eager greedy search reaches 1.05 seconds and
the no-adaptation baseline remains near 0.10 seconds. Plan lengths increase for both
search strategies and remain similar, reaching 46.00 and 46.80 actions, respectively.
Shaded bands indicate standard deviations.}
\end{figure}

\begin{figure}[tb]
    \centering
    \begin{subfigure}[t]{0.44\linewidth}
        \centering
        % Generated by navigation_planta; do not edit by hand.
\begin{tikzpicture}[baseline=(paperAxis.north)]
\begin{axis}[name=paperAxis,
width=\textwidth,
height=0.75\textwidth,
xlabel={Number of nodes},
ylabel={Planning Time [seconds]},
xmin=9,
xmax=90,
ymin=0,
ymax=220,
xtick={10,20,30,40,50,60,70,80,90},
grid=major,
legend pos=north west,
legend style={font=\footnotesize,nodes={scale=0.8,transform shape}},
legend image post style={scale=0.5}
]
\addplot[name path=upper0,draw=none,forget plot] table[x=nodes,y=mean_time_plus_std,col sep=comma] {pddl/navigation_experiment/data/s5_astar_experiment_processed.csv};
\addplot[name path=lower0,draw=none,forget plot] table[x=nodes,y=mean_time_minus_std,col sep=comma] {pddl/navigation_experiment/data/s5_astar_experiment_processed.csv};
\addplot[fill=\greenColor!20,fill opacity=0.3,draw=none,forget plot] fill between[of=upper0 and lower0];
\addplot[thick,color=\greenColor] table[x=nodes,y=mean_time,col sep=comma] {pddl/navigation_experiment/data/s5_astar_experiment_processed.csv};
\addlegendentry{A* mean}
\addplot[name path=upper1,draw=none,forget plot] table[x=nodes,y=mean_time_plus_std,col sep=comma] {pddl/navigation_experiment/data/s5_eager_experiment_processed.csv};
\addplot[name path=lower1,draw=none,forget plot] table[x=nodes,y=mean_time_minus_std,col sep=comma] {pddl/navigation_experiment/data/s5_eager_experiment_processed.csv};
\addplot[fill=\blueColor!20,fill opacity=0.3,draw=none,forget plot] fill between[of=upper1 and lower1];
\addplot[thick,color=\blueColor] table[x=nodes,y=mean_time,col sep=comma] {pddl/navigation_experiment/data/s5_eager_experiment_processed.csv};
\addlegendentry{EG mean}
\end{axis}
\end{tikzpicture}

        \caption{Planning time as map size increases.}%
        \label{planta:S-plot:s5_planning_time}
    \end{subfigure}
    \begin{subfigure}[t]{0.44\linewidth}
        \centering
        % Generated by navigation_planta; do not edit by hand.
\begin{tikzpicture}[baseline=(paperAxis.north)]
\begin{axis}[name=paperAxis,
width=\textwidth,
height=0.75\textwidth,
xlabel={Number of nodes},
ylabel={Plan Length [no. of actions]},
xmin=9,
xmax=90,
ymin=0,
ymax=65,
xtick={10,20,30,40,50,60,70,80,90},
grid=major,
legend pos=north west,
legend style={font=\footnotesize,nodes={scale=0.8,transform shape}},
legend image post style={scale=0.5}
]
\addplot[name path=upper0,draw=none,forget plot] table[x=nodes,y=mean_action_count_plus_std,col sep=comma] {pddl/navigation_experiment/data/s5_astar_experiment_processed.csv};
\addplot[name path=lower0,draw=none,forget plot] table[x=nodes,y=mean_action_count_minus_std,col sep=comma] {pddl/navigation_experiment/data/s5_astar_experiment_processed.csv};
\addplot[fill=\greenColor!20,fill opacity=0.3,draw=none,forget plot] fill between[of=upper0 and lower0];
\addplot[thick,color=\greenColor] table[x=nodes,y=mean_action_count,col sep=comma] {pddl/navigation_experiment/data/s5_astar_experiment_processed.csv};
\addlegendentry{A* mean}
\addplot[name path=upper1,draw=none,forget plot] table[x=nodes,y=mean_action_count_plus_std,col sep=comma] {pddl/navigation_experiment/data/s5_eager_experiment_processed.csv};
\addplot[name path=lower1,draw=none,forget plot] table[x=nodes,y=mean_action_count_minus_std,col sep=comma] {pddl/navigation_experiment/data/s5_eager_experiment_processed.csv};
\addplot[fill=\blueColor!20,fill opacity=0.3,draw=none,forget plot] fill between[of=upper1 and lower1];
\addplot[thick,color=\blueColor] table[x=nodes,y=mean_action_count,col sep=comma] {pddl/navigation_experiment/data/s5_eager_experiment_processed.csv};
\addlegendentry{EG mean}
\end{axis}
\end{tikzpicture}

        \caption{Plan length as map size increases.}%
        \label{planta:S-plot:s5_action_number}
    \end{subfigure}
    \caption{\Sfive results: Planning time and plan length as map size increases. The shaded area shows the standard deviation. EG stands for the eager greedy search.}%
    \label{planta:S-fig:s5_results}
\Description{Two line-plot panels show mean planning time and mean plan length against map size, from
9 to 90 retained nodes, in the combined stress test. A* planning time rises to 188.96
seconds, whereas eager greedy search reaches 13.73 seconds. Plan lengths increase with
map size, reaching 48.10 actions for A* and 51.60 for eager greedy search. Shaded bands
indicate standard deviations.}
\end{figure}
